% Part: normal-modal-logic
% Chapter: filtrations
% Section: introduction

\documentclass[../../../include/open-logic-section]{subfiles}

\begin{document}

\olfileid{nml}{fil}{int}

\olsection{Pambuka}

Pitakon wigati ngenani sawijining logika yaiku apa logika iku bisa
diputusake, yaiku apa ana prosedur efektif kang bisa mangsuli pitakon
``apa !!{formula} iki valid?'' Logika proposisional bisa diputusake:
kita bisa nguji kanthi efektif apa !!a{formula} iku tautologi kanthi
mbangun tabel bebener, lan kanggo !!{formula} tartamtu tabel kasebut
winates. Nanging, ora langsung cetha kepiye nguji apa sawijining
formula modal bener ing kabeh model, amarga cacah modele tanpa wates.
Kita bisa ndhaptar kabeh model winates kang relevan kanggo sawijining
formula, amarga mung panetepan himpunan bagean jagad marang
!!{propositional variable}s kang pancen katon ing !!{formula} iku
relevan. Yen relasi aksesibilitas wis tetep, saben panetepan $V(p)$
sing mungkin mung sawijining himpunan bagean saka~$W$, lan yen
$\card{W} = n$ ana $2^n$ pilihan. Yen !!{formula}~$!A$ ngemot $m$
!!{propositional
  variable}s, ana $2^{nm}$ model beda kanthi~$n$ jagad. Kanggo saben
model, kita bisa nguji apa~$!A$ bener ing kabeh jagad kanthi ngetung
nilai bebener~$!A$ ing saben jagad. Mesthi wae, kita uga kudu mriksa
kabeh relasi aksesibilitas kang mungkin, nanging cacah relasi ing~$n$
jagad uga winates (persise, cacah himpunan bagean saka $W \times W$,
yaiku $2^{n^2}$).

Yen kang kita teliti dudu logika~\Log{K}, nanging logika kang
ditepakake dening sawijining kelas model (upamane model refleksif lan
transitif), kita uga kudu bisa nguji apa relasi aksesibilitas iku saka
jinis kang bener. Iki bisa ditindakake yen bingkai kang kita teliti
bisa didéfinisikake dening !!{formula}s modal (upamane kanthi nguji apa
\Ax{T} lan \Ax{4} valid ing bingkai kasebut). Dadi, kita bisa nlusuri
kabeh bingkai winates, nguji apa saben bingkai kalebu kelas kang
diteliti, banjur ndhaptar kabeh model kang mungkin ing bingkai iku
lan nguji apa $!A$ bener ing saben model. Yen ora, mandhega:
$!A$ ora valid ing kelas model kang diteliti.

Ana masalah karo gagasan iki: kita ora ngerti kapan, yen pancen ana,
panelusuran bisa dipungkasi. Yen !!{formula} nduweni kontramodel
winates, prosedur kita bakal nemokake. Nanging yen ora nduweni
kontramodel winates, kita ora bakal entuk jawaban. !!{formula} iku bisa
wae valid (ora ana kontramodel babar pisan), utawa mung nduweni
kontramodel tanpa wates kang ora bakal kita priksa. Masalah iki bisa
dirampungi yen kita bisa nuduhake yen saben !!{formula} kang nduweni
kontramodel uga nduweni kontramodel winates. Yen mangkono, kita kandha
yen logika iku nduweni \emph{sipat model winates}.

Nanging kepiye mbuktekake yen sawijining logika nduweni sipat model
winates? Salah siji cara yaiku nemokake cara ngowahi (kontra)model
tanpa wates saka~$!A$ dadi model winates. Yen iki bisa ditindakake,
saben ana model kang $!A$ ora bener, model winates kang diasilake uga
ndadekake $!A$ ora bener. Model winates iku bakal katon ing dhaptar
kabeh model winates, lan pungkasane kita bisa nemtokake kanggo saben
formula kang ora valid yen pancen ora valid. Prosedur iki ora bakal
mandheg yen formulane valid. Yen saliyane iku kita bisa nuduhake yen
ana ukuran maksimum kanggo model winates kang diwenehake prosedur,
lan yen ukuran maksimum iku mung gumantung marang
!!{formula}~$!A$, kita bakal duwe wates ukuran kanggo nguji model
winates nalika nggoleki kontramodel. Yen nganti wates kasebut kita
durung nemokake kontramodel, tegesé pancen ora ana. Mula prosedur kita
bakal bisa mutusake pitakon ``apa $!A$ valid?'' kanggo sembarang
!!{formula}~$!A$.

Strategi kang kerep bisa ngowahi struktur tanpa wates dadi struktur
winates yaiku ``ngidhentifikasi'' !!{element}s saka struktur kang
prilakune padha tumrap perkara kang relevan. Yen ana tanpa wates
jagad ing~$\mModel{M}$ kang prilakune padha tumrap perkara kang
relevan, bisa uga kita nglumpukake \emph{kabeh} jagad iku dadi
sawetara ``kelas'' kang cacahé winates (senajan saben kelas bisa tanpa
wates). Kanthi tembung liya, kita kudu mbagi himpunan jagad kanthi
cara kang trep, yaiku saben kelas bisa ngemot tanpa wates jagad,
nanging cacah kelase winates. Banjur kita netepake
model anyar~$\mModel{M^*}$ kang jagade yaiku kelas-kelas kasebut.
Kelas kang cacahé winates ing model lawas menehi jagad kang cacahé
winates ing model anyar, yaiku model winates. Sebuta kelas kang ngemot
jagad~$w$ kanthi $[w]$. Kita kepengin $\mSat{M}{!A}[w]$ yen lan mung
yen $\mSat{M^*}{!A}[{[w]}]$, amarga kita kepengin model anyar dadi
kontramodel tumrap~$!A$ yen model lawas uga mangkono. Kanggo iki,
kita kudu netepake pambagian kelas lan relasi aksesibilitas
saka~$\mModel{M^*}$ kanthi trep.

Kanggo ndeleng cara kerjane, upamane dhisik relasi aksesibilitas iku
universal. $\mSat{M}{\Box !B}[w]$ yen lan mung yen kanggo saben
$v \in W$, $\mSat{M}{!B}[v]$; padha uga kanggo $\mModel{M^*}$,
nanging nganggo $[w]$ lan $[v]$. Minangka gagasan wiwitan, anggepa
loro jagad $u$ lan $v$ ekuivalen (kalebu kelas kang padha) yen loro-lorone
sarujuk tumrap kabeh !!{propositional variable}s ing~$\mModel{M}$,
yaiku $\mSat{M}{p}[u]$ yen lan mung yen $\mSat{M}{p}[v]$.
Tetepa $V^*(p) = \Setabs{[w]}{\mSat{M}{p}[w]}$. Tujuwan kita yaiku
nuduhake $\mSat{M}{!A}[w]$ yen lan mung yen
$\mSat{M^*}{!A}[{[w]}]$. Mesthi iki dibuktekake kanthi induksi:
kasus dhasar yaiku $!A \equiv p$. Sepisan, upamane
$\mSat{M}{p}[w]$. Banjur $[w] \in V^*(p)$ miturut definisi, mula
$\mSat{M^*}{p}[{[w]}]$. Saiki upamane
$\mSat{M^*}{p}[{[w]}]$. Iki ateges $[w] \in V^*(p)$, yaiku kanggo
sawetara $v$ kang ekuivalen karo~$w$, $\mSat{M}{p}[v]$. Nanging
``$w$ ekuivalen karo $v$'' ateges ``$w$ lan $v$ ndadekake kabeh
!!{propositional variable}s kang padha bener,'' mula
$\mSat{M}{p}[w]$. Saiki kanggo langkah induksi, upamane $!A
\ident \lnot !B$. Banjur $\mSat{M}{\lnot !B}[w]$ yen lan mung yen
$\mSat/{M}{!B}[w]$; yen lan mung yen
$\mSat/{M^*}{!B}[{[w]}]$ (miturut hipotesis induksi); yen lan mung yen
$\mSat{M^*}{\lnot !B}[{[w]}]$. Padha uga tumrap operator nonmodal
liyane. Iki uga lumaku kanggo $\Box$: upamane
$\mSat{M^*}{\Box
  !B}[{[w]}]$. Iki ateges kanggo saben $[u]$,
$\mSat{M^*}{!B}[{[u]}]$. Miturut hipotesis induksi, kanggo saben $u$,
$\mSat{M}{!B}[u]$. Mula, $\mSat{M}{\Box !B}[w]$.

Ing kasus umum, nalika kita uga kudu netepake relasi aksesibilitas
kanggo $\mModel{M^*}$, perkarane luwih rumit. Model~$\mModel{M^*}$
bakal kita sebut \emph{filtrasi} yen relasi aksesibilitase~$R^*$
nyukupi sarat kang dibutuhake supaya bukti induksi ing ndhuwur lumaku.
Banjur saben filtrasi~$\mModel{M^*}$ ndadekake $!A$ bener ing $[w]$
yen lan mung yen $\mModel{M}$ ndadekake $!A$ bener ing~$w$. Nanging
saiki kita uga kudu nuduhake yen filtrasi \emph{ana}, yaiku kita bisa
netepake~$R^*$ saengga nyukupi sarat kang dibutuhake. Supaya cara iki
bisa lumaku, jagad $u$ lan $v$ kudu dianggep ekuivalen ora mung yen
padha sarujuk tumrap kabeh !!{propositional variable}s, nanging uga
tumrap kabeh sub-!!{formula}s saka~$!A$. Amarga $!A$ mung nduweni
sub-!!{formula}s kang cacahé winates, filtrasi uga dijamin winates.
Filtrasi bisa ditepakake kanthi luwih saka siji cara. Supaya relasi
aksesibilitas filtrasi nyukupi sipat kang dibutuhake (upamane
refleksif, transitif, lan sapanunggalane), kita kudu milih definisi
kanggo~$R^*$ kanthi ngati-ati.


\end{document}


